\section{Điều các họ chỉ số chưa giải quyết}
\label{sec:ch08-dieu-chua-giai-quyet}

\index{độ trung thực!kiểm định chỉ số}%
\index{chỉ số đánh giá!bốn nhóm}%
Bốn nhóm ở các mục trên được xếp theo ba ô của
hình~\ref{fig:ch08-vong-lap-do}, tức theo cách sách này đọc chúng. Bài báo neo
của cuốn sách cũng xếp các chỉ số thành nhóm, độc lập với cách xếp ấy, và kết
quả gần trùng nhau đủ để đáng ghi lại. Gur-Arieh và cộng sự chia các chỉ số họ
đánh giá thành bốn nhóm theo đúng câu hỏi \enquote{chỉ số này can thiệp vào
cái gì}: nhóm dựa trên mức quan trọng, nhóm dựa trên tham số, nhóm dựa trên
attribution, và nhóm dựa trên
\tn{tính hữu dụng ngữ nghĩa}{semantic utility}~\autocite{p21metrics}.

Trước khi dùng cách chia ấy, phải nói rõ một điều về phạm vi, vì bỏ qua nó sẽ
làm hỏng cách đọc cả Phần~IV. Đối tượng của bài báo neo không phải attribution
đặc trưng, mà là \tn{chuỗi suy luận}{chain-of-thought} của
\tn{mô hình ngôn ngữ lớn}{large language model}:
nó hỏi lời kể lại các bước suy luận có phản ánh đúng phép tính đã dẫn tới câu
trả lời hay không. Không tên chỉ số nào của họ xóa dần và họ tiên đề xuất hiện trong bài báo đó, và
ngược lại, tám chỉ số mà nó đánh giá đều được thiết kế riêng cho chuỗi suy
luận~\autocite{p21metrics}. Hai bên đo hai đối tượng khác nhau bằng hai bộ
công cụ khác nhau.

Cái đi qua được ranh giới ấy là cấu trúc, không phải kết quả. Nhóm dựa trên mức
quan trọng làm đúng việc mà mục~\ref{sec:ch08-ho-xoa-dan} đã mô tả: nó nhiễu một
phần của lời giải thích rồi xem câu trả lời có đổi không. Nhóm dựa trên tham số
can thiệp vào trọng số của mô hình thay vì vào đầu vào. Trong bốn mục trên, chỉ
phép thử thứ nhất ở mục~\ref{sec:ch08-dieu-kien-can} động tới trọng số, mà nó
là điều kiện cần chứ không phải một chỉ số, nên ở phía attribution nhóm này
chưa có đại diện thực sự. Nhóm dựa trên attribution dùng lại chính SHAP của
chương~\ref{ch:shap} làm dụng cụ đo. Nhóm thứ tư thì đo thứ khác hẳn: lời giải
thích có đủ dùng cho một mô hình yếu hơn tái lập được câu trả lời hay
không~\autocite{p21metrics}.

Chẩn đoán mà bài báo đưa ra cho hai nhóm lớn nhất là chỗ đáng mang về nhất, vì
nó nói đúng cái mà chương này đi tìm: giả định nào đã lẻn vào. Với nhóm dựa trên
mức quan trọng, các tác giả cho rằng chúng đồng nhất mức quan trọng với độ trung
thực, nên có thể gán nhãn không trung thực cho một bước vốn trung thực mà chỉ vì
nó không thiết yếu với câu trả lời cuối~\autocite{p21metrics}. Đó chính là phép
đồng nhất mà cả họ xóa dần đứng trên: bỏ đi một đặc trưng mà đầu ra không đổi
được đọc thành đặc trưng ấy không được mô hình dùng, trong khi nó chỉ có nghĩa
là đặc trưng ấy không quyết định một mình. Với nhóm dựa trên tính hữu dụng ngữ
nghĩa, chẩn đoán ngược lại: một lời biện minh nghe hợp lý nhưng bịa đặt vẫn đủ
để một mô hình yếu hơn tái lập câu trả lời~\autocite{p21metrics}. Đây là sự lẫn
lộn giữa độ trung thực và \tn{tính hợp lý}{plausibility} mà
mục~\ref{sec:ch01-do-trung-thuc-va-tinh-hop-ly} đã tách ra, lần này nằm ngay
trong dụng cụ đo.

Còn lại một câu hỏi mà không mục nào trong chương trả lời được, và nó là câu hỏi
của cả Phần~IV. Mỗi họ chỉ số ở trên được đề xuất kèm một lập luận vì sao con số
của nó đáng tin, và không họ nào được đối chiếu với ground truth lấy từ bên
ngoài chính nó. Bài báo neo nói thẳng tình trạng ấy: phần lớn công trình đề xuất
chỉ số chỉ báo cáo điểm số tuyệt đối hoặc so với các chỉ số trước, còn số ít bộ
đánh giá đã có thì dựa vào những đại lượng thay thế như tính hợp lý hoặc mức
quan trọng, vốn là các tính chất trực giao với độ trung thực~\autocite{p21metrics}.

Nói theo phép so sánh của mục~\ref{sec:ch08-chi-so-nhu-mot-dung-cu}: cả một tủ
dụng cụ đo đã được chế tạo, được hiệu chỉnh lẫn nhau, và được dùng để xếp hạng
các phương pháp, mà chưa cái nào được đưa ra đối chiếu với ground truth. Các
tác giả của bài báo neo cũng nêu lý do, và lý do ấy không phải sự cẩu thả:
ground truth ở đây đòi biết mô hình thật sự đã dùng gì, mà nội dung tính toán
bên trong mô hình thì không quan sát trực tiếp được~\autocite{p21metrics}.

Muốn thoát khỏi vòng ấy thì phải dựng lấy ground truth, bằng cách chọn những
bài toán mà đầu ra của chúng để lộ phép tính đã sinh ra nó.
Chương~\ref{ch:chi-so-khong-do-do-trung-thuc} đọc kỹ công trình đã làm đúng việc
đó, và đọc luôn kết quả mà nó thu được khi đem các chỉ số ra đối chiếu.
